\documentclass[11pt,a4paper]{article}
\usepackage{amsmath}
\usepackage{amssymb}
\usepackage{graphicx}
\usepackage{fontspec}
\usepackage{babel}
\usepackage[margin=2.5cm]{geometry}
\usepackage{enumitem}
\usepackage{booktabs}
\usepackage{xcolor}
\usepackage{hyperref}

\definecolor{coral}{HTML}{cc785c}
\definecolor{ink}{HTML}{141413}
\definecolor{muted}{HTML}{6c6a64}

\hypersetup{colorlinks=true, linkcolor=coral, urlcolor=coral}

\babelprovide[import, onchar=ids fonts]{thai}
\babelprovide[import, onchar=ids fonts]{english}

\babelfont{rm}{Noto Sans}
\babelfont[thai]{rm}{Noto Sans Thai}
\babelfont{tt}{Noto Sans Mono}
\babelfont[thai]{tt}{Noto Sans Thai}

\directlua{
  local glyph_id = node.id("glyph")
  local penalty_id = node.id("penalty")

  local function is_thai(c)
    return c >= 0x0E01 and c <= 0x0E5B
  end

  local function is_thai_combining(c)
    return (c >= 0x0E31 and c <= 0x0E3A) or (c >= 0x0E47 and c <= 0x0E4E)
  end

  luatexbase.add_to_callback("pre_linebreak_filter", function(head)
    for n in node.traverse_id(glyph_id, head) do
      local prev = n.prev
      if prev and prev.id == glyph_id and is_thai(prev.char)
         and is_thai(n.char) and not is_thai_combining(n.char) then
        local p = node.new(penalty_id)
        p.penalty = 0
        node.insert_before(head, n, p)
      end
    end
    return head
  end, "thai_break")
}

\emergencystretch=1em

\begin{document}

\begin{center}
  {\LARGE \textbf{สรุปงานวิจัยที่เกี่ยวข้อง}}\\[6pt]
  {\large การสร้างคำอธิบายรูปภาพและกราฟในเอกสารวิชาการ}\\[4pt]
  {\color{muted}\small Figure \& Chart Captioning in Scientific Documents}
\end{center}

\noindent\rule{\textwidth}{1.5pt}

\vspace{16pt}

%% ─────────────────────────────────────────────────────────────────
\section*{ภาพรวม (Overview)}

งานวิจัยที่รวบรวมไว้ครอบคลุมพัฒนาการของการสร้างคำอธิบายรูปภาพและกราฟในเอกสารวิชาการ
ตั้งแต่การสร้างชุดข้อมูล (2021) จนถึงการใช้ Multimodal LLM แบบร่วมมือ (2025)
โดยสามารถจัดกลุ่มได้เป็น 3 แนวทางหลัก:
\textbf{(1)} การสร้างชุดข้อมูลขนาดใหญ่,
\textbf{(2)} การแปลง plot เป็นข้อความเพื่อ reasoning, และ
\textbf{(3)} การใช้ LLM แบบหลายตัวร่วมกัน

\vspace{12pt}
\noindent\rule{\textwidth}{0.5pt}

%% ─────────────────────────────────────────────────────────────────
\section*{Paper 1 — SciCap (2021)}

\noindent{\color{coral}\textbf{ชื่อ:}} Generating Captions for Scientific Figures\\
\noindent{\color{muted}\small Hsu, Giles, Huang · Pennsylvania State University · arXiv:2110.11624}

\vspace{8pt}
\setlength{\parindent}{2em}
งานนี้เสนอ \textbf{SciCap} ซึ่งเป็นชุดข้อมูลขนาดใหญ่สำหรับการสร้างคำอธิบาย
รูปภาพในบทความวิชาการ โดยดึงข้อมูลจากบทความ arXiv สาขา Computer Science
ที่ตีพิมพ์ระหว่างปี 2010--2020 รวมกว่า 290,000 บทความ
ผ่านกระบวนการ pre-processing ได้แก่
การจำแนกประเภทรูปภาพ, การระบุ sub-figure, การ normalize ข้อความ
และการเลือก caption text ทำให้ได้รูปภาพกว่า \textbf{2 ล้านรูป}

\vspace{6pt}
\noindent\textbf{จุดเด่น:}
\begin{itemize}[nosep, leftmargin=2em]
  \item ชุดข้อมูลแรกในระดับ large-scale สำหรับ scientific figure captioning
  \item มีระบบจำแนก figure type — graph plot เป็น type หลัก (19.2\%)
  \item สร้าง baseline model สำหรับ graph plot captioning
  \item เปิดเผยทั้งโอกาสและความท้าทายของงานนี้
\end{itemize}

\vspace{12pt}
\noindent\rule{\textwidth}{0.4pt}

%% ─────────────────────────────────────────────────────────────────
\section*{Paper 2 — SciCap+ (2023)}

\noindent{\color{coral}\textbf{ชื่อ:}} A Knowledge Augmented Dataset to Study the Challenges of Scientific Figure Captioning\\
\noindent{\color{muted}\small arXiv:2306.03491}

\vspace{8pt}
\setlength{\parindent}{2em}
\textbf{SciCap+} ต่อยอดจาก SciCap โดยเพิ่ม \textbf{knowledge augmentation}
เข้าไปในชุดข้อมูล เพื่อศึกษาความท้าทายของการสร้างคำอธิบาย
รูปภาพวิทยาศาสตร์ให้ลึกขึ้น งานนี้ขยายขอบเขตการเข้าใจเอกสารวิชาการ
ออกไปเกินกว่าการวิเคราะห์ข้อความล้วน โดยรวมบริบทจากเนื้อหาของบทความเข้ามาด้วย

\vspace{6pt}
\noindent\textbf{จุดเด่น:}
\begin{itemize}[nosep, leftmargin=2em]
  \item เพิ่ม knowledge context จากเนื้อหาบทความให้แต่ละรูปภาพ
  \item ศึกษาความท้าทายที่ชุดข้อมูลดั้งเดิม (SciCap) ยังไม่ครอบคลุม
  \item เหมาะสำหรับการพัฒนา context-aware captioning
\end{itemize}

\vspace{12pt}
\noindent\rule{\textwidth}{0.4pt}

%% ─────────────────────────────────────────────────────────────────
\section*{Paper 3 — DePlot (2023)}

\noindent{\color{coral}\textbf{ชื่อ:}} One-shot Visual Language Reasoning by Plot-to-Table Translation\\
\noindent{\color{muted}\small Liu et al. · Google DeepMind · arXiv:2212.10505}

\vspace{8pt}
\setlength{\parindent}{2em}
\textbf{DePlot} นำเสนอแนวทาง few-shot สำหรับ visual language reasoning
โดยแยกปัญหาออกเป็น 2 ขั้นตอน: \textbf{(1)} แปลง plot/chart เป็นตารางข้อความ
(plot-to-table translation) และ \textbf{(2)} ใช้ LLM reasoning บนข้อความที่ได้
โมดูล DePlot ทำหน้าที่แปลง modality โดยเฉพาะ
แล้วส่งผลลัพธ์ไปยัง LLM ใด ๆ แบบ plug-and-play

\vspace{6pt}
\noindent\textbf{ผลลัพธ์:}
\begin{itemize}[nosep, leftmargin=2em]
  \item DePlot + LLM (one-shot) ดีกว่า SOTA ที่ fine-tune บนข้อมูลหลายพัน samples
  \item \textbf{+29.4\%} บน ChartQA (human-written queries)
  \item ไม่ต้องการ fine-tuning บน downstream task — ใช้ได้ทันที
\end{itemize}

\vspace{12pt}
\noindent\rule{\textwidth}{0.4pt}

%% ─────────────────────────────────────────────────────────────────
\section*{Paper 4 — MLBCAP (2025)}

\noindent{\color{coral}\textbf{ชื่อ:}} Multi-LLM Collaborative Caption Generation in Scientific Documents\\
\noindent{\color{muted}\small Kim, Lee, Choi, Hsu et al. · Teamreboott / Penn State / Adobe Research · arXiv:2501.02552}

\vspace{8pt}
\setlength{\parindent}{2em}
\textbf{MLBCAP} เสนอ framework ที่ใช้ LLM หลายตัวทำงานร่วมกัน
เพื่อสร้าง caption สำหรับรูปภาพในเอกสารวิชาการ
แก้ปัญหาของวิธีการเดิมที่ใช้ข้อมูลไม่ครบถ้วน
และมองงาน captioning แบบ isolated โดยไม่คำนึงถึงบริบทรอบข้าง

\vspace{6pt}
\noindent\textbf{จุดเด่น:}
\begin{itemize}[nosep, leftmargin=2em]
  \item ใช้ LLM หลายตัวร่วมมือกัน (collaborative multi-LLM)
  \item รวมข้อมูลที่หลากหลายเข้าด้วยกัน ไม่ใช่แค่ภาพเดียว
  \item ต่อยอดจากงานของ SciCap series (Hsu et al.)
\end{itemize}

\vspace{12pt}
\noindent\rule{\textwidth}{0.4pt}

%% ─────────────────────────────────────────────────────────────────
\section*{Paper 5 — Survey: Chart Understanding with MLLMs (2025)}

\noindent{\color{coral}\textbf{ชื่อ:}} Multimodal Information Fusion for Chart Understanding: A Survey of MLLMs\\
\noindent{\color{muted}\small arXiv:2602.10138}

\vspace{8pt}
\setlength{\parindent}{2em}
งาน survey นี้รวบรวมพัฒนาการของ Multimodal LLM (MLLM) สำหรับ
chart understanding โดยครอบคลุมการผสานข้อมูลหลาย modality
(กราฟิกและข้อความ) เพื่อดึงความหมาย พร้อมวิเคราะห์ข้อจำกัด
และแนวทางการปรับปรุงด้วย cognitive enhancement

\vspace{6pt}
\noindent\textbf{ประเด็นที่ครอบคลุม:}
\begin{itemize}[nosep, leftmargin=2em]
  \item วิวัฒนาการของ MLLM สำหรับ chart understanding
  \item ข้อจำกัดของโมเดลปัจจุบัน
  \item แนวทาง cognitive enhancement เพื่อเพิ่มความสามารถในการ reasoning
\end{itemize}

\vspace{16pt}
\noindent\rule{\textwidth}{1.5pt}

%% ─────────────────────────────────────────────────────────────────
\section*{ตารางเปรียบเทียบ}

\vspace{8pt}
\noindent
\begin{tabular}{@{}lllp{5.5cm}@{}}
  \toprule
  \textbf{Paper} & \textbf{ปี} & \textbf{แนวทาง} & \textbf{จุดเด่น} \\
  \midrule
  SciCap     & 2021 & Dataset       & 2M+ figures จาก 290K papers \\
  SciCap+    & 2023 & Dataset+      & Knowledge augmentation \\
  DePlot     & 2023 & Plot→Table→LLM & +29.4\% บน ChartQA (one-shot) \\
  MLBCAP     & 2025 & Multi-LLM     & Collaborative captioning \\
  Survey     & 2025 & Survey        & MLLM evolution \& limitations \\
  \bottomrule
\end{tabular}

\vspace{16pt}
\noindent\rule{\textwidth}{1.5pt}
\vspace{6pt}

\noindent{\color{muted}\small สรุปโดย Claude Sonnet 4.6 · \today}

\end{document}
